\documentclass[11pt,a4paper]{article}
\usepackage[margin=25mm]{geometry}
\usepackage{amsmath,amssymb,amsthm,booktabs,hyperref}
\newtheorem{theorem}{Theorem}
\newtheorem{lemma}{Lemma}
\newcommand{\gen}[1]{\langle#1\rangle}
\title{Small primitive strings with intransitive omissions}
\author{Human Compression project: bounded L044 component}
\date{}
\begin{document}\maketitle
We recover Proposition~5.10 and the sixteen literal coloured graphs of
Table~2 in Cameron--Fernandes--Leemans--Mixer,
\emph{Highest rank of a polytope for $A_n$}, arXiv:1605.09173v1,
physical pages~26--27. The source attributes this finite classification
to a Magma computation. Our replacement uses the previously accepted
prefix method with a new necessary orbit test and an exact correspondence
with the source drawings. It is a required finite proof component,
not a diagnostic. Proposition~5.11's subsequent attachment exclusions,
Table~1's different rank maxima, and the full original remain outside
this note.

\begin{theorem}\label{thm:small}
Let $G\leq A_n$ be primitive on $n$ points, $2\leq n\leq11$,
with a string C-marking $(g_0,\ldots,g_{d-1})$ of nonidentity involutions
and connected string diagram. Suppose every maximal marked omission
is intransitive and $d\geq(n-1)/2$. Up to simultaneous permutation
conjugacy and reversal of the entire string, its coloured permutation
graph is exactly one of the sixteen graphs in L044 Table~2.
There are 27 classes before taking reversal.
\end{theorem}

\begin{center}
\begin{tabular}{lrrrr}
\toprule
Source rows & degree & rank & ordered classes & reversal classes\\
\midrule
1: $D_{10}$ &5&2&1&1\\
2--3: $A_5\cong\mathrm{PSL}_2(5)$ &6&3&3&2\\
4--9: $A_9$ &9&4&11&6\\
10--13: $A_{10}$ &10&5&8&4\\
14--16: $A_{11}$ &11&6&4&3\\
\bottomrule
\end{tabular}
\end{center}
Here an ordered class fixes every colour; only the separately stated
full reversal identifies dual strings. No arbitrary reordering is used.
The actual permutation isomorphism $A_5\to\mathrm{PSL}_2(5)$ is the
classical retained naming convention; the finite program identifies
the degree-six group with $A_5$ explicitly.

\begin{lemma}[Complete prefix enumeration]\label{lem:engine}
The programs \texttt{census.g} and \texttt{run.g} exhaust all markings
in Theorem~\ref{thm:small}, modulo simultaneous conjugation.
\end{lemma}
\begin{proof}
We retain the completeness of GAP's primitive-group catalogue in
degrees $2$ through $11$. Every action is visited, including those
with no output. A group containing an odd permutation is ineligible.
If its involutions do not generate it, no involutory marking generates
it. These exact tests leave seventeen actions.

For each remaining $G$, begin with the empty tuple. At a prefix
$\mathbf g$ with group $H$, consider every involution outside $H$
commuting with all but the last prefix generator and not commuting
with the last. These are precisely the necessary next-letter
conditions for an independent connected string. Conjugating by the
stabilizer of the whole prefix in $N_{S_n}(G)$ preserves every
condition; retaining one representative of each such orbit therefore
loses no conjugacy class. The next stabilizer is its centralizer of
the chosen involution.

Every accepted prefix is checked for independence. It must also
satisfy
\[
 \gen{g_0,\ldots,g_{l-1}}\cap\gen{g_1,\ldots,g_l}
       =\gen{g_1,\ldots,g_{l-1}}.
\]
This is used only as a necessary condition; no assertion that this
single test proves all IP identities is needed. If a proper prefix
is already transitive, then in any longer marking the omission of
the final generator contains that transitive prefix, which is
forbidden. Such a prefix can be discarded.

Put $t=\lfloor n/2\rfloor$. A prefix of length $l$ needs at least
$t-l$ further independent letters, so its index must be at least
$2^{\max(0,t-l)}$. If $l>0$, the letters after the next one generate
a subgroup $T\leq C_G(H)$ of order at least
$2^{q}$, $q=\max(0,t-l-1)$. Full IP in any eventual valid marking
gives $T\cap H=1$, hence $T\cap Z(H)=1$ and
\[
                     |C_G(H):Z(H)|\geq2^q.
\]
Failure of either integer bound safely discards the prefix.
The reused Whiston independent-family bound for a proper subgroup
of $S_n$ gives length at most $n-2$; thus this is the complete
upper search limit, not a rank guessed from the source table.

When $H=G$, the program checks the target rank and every maximal
omission's orbit partition. It then constructs every marked
parabolic and checks all $2^{2d}$ intersection equalities.
Consequently every output is an eligible string C-marking, and
the preceding necessary conditions ensure completeness.
\end{proof}

\begin{proof}[Proof of Theorem~\ref{thm:small}]
The complete run gives the 27 ordered classes in the displayed
table and no others. Their permutation tuples are written to the
fresh correspondence receipt. The compulsory input
\texttt{source\_table2.json} independently transcribes every edge
of each of the sixteen source drawings, including parallel edges
with different colours; the coloured matchings reconstruct the
literal involutions. Source vertices are numbered left to right
on the upper path and then left to right below it.

For two such transitive coloured graphs, an isomorphism is uniquely
determined by the image of any one vertex: propagate it along
each colour, retaining loops at fixed points. The checker tries
every possible image, rejects inconsistent or noninjective maps,
and accepts only a complete bijection. Thus it checks actual
simultaneous permutation conjugacy, not agreement of orders or
Schl\"afli types. It does this both in the stated order and after
reversing the full source tuple. Every computed marking matches
one source row, every source row occurs, distinct source rows are
inequivalent even under reversal, and the 27 computed ordered
classes are pairwise inequivalent without reversal.

The degree-five group is generated by two involutions whose product
has order five, and is $D_{10}$. GAP constructs an isomorphism of
the degree-six group with $A_5$. The larger output groups are
checked as the actual $A_9,A_{10},A_{11}$ in their natural actions.
This verifies the group column as well as all literal markings.
\end{proof}

\paragraph{Proof and software boundary.}
The complete compulsory finite component consists of
\texttt{census.g}, \texttt{run.g}, \texttt{check.py}, and
\texttt{source\_table2.json}. The Python checker performs the
mathematical coloured-graph correspondence as well as launching
GAP. GAP permutation-group arithmetic and its PrimGrp catalogue
are retained software/classification inputs, with full unknown
external proof cost. The earlier L047/common theory supplies
Whiston's stated bound and the already reviewed prefix-method
argument; these are paid dependencies, not a new theorem of this
note. No earlier census output is used instead of the fresh run.
The source page was directly viewed for all sixteen drawings.
There is no new formalization, original admission, question
solution, or priority claim.
\end{document}
